% ===== 原文 pp.173–187（PDF p.188–202）· tr-05 · c05a =====
% 本块 = 第 5 章标题 + §5.1（含 5.1.1、5.1.2 两个三级标题）+ §5.2 + §5.3 开头；
% tag 5.1–5.4；图 Figure 1–5 已裁为 work/figures/f5_1.png … f5_5.png。
% 块末按 STYLE_GUIDE §10「跨页句子开始者写完」补完 p.187 末尾跨到 p.188 的那一句。

\section{周期积分}

% [原文 p.173]
在第 4 章中，我们证明了环面曲面中曲线计数的 Mikhalkin 公式。特殊化到 $\PP^2$ 的情形，基本上就得到 $\PP^2$ 的 A-模型，因为这些公式计算 Gromov--Witten 不变量 $\langle T_2^{3d-1}\rangle_{0,d}$，即穿过 $3d-1$ 个点的 $d$ 次有理曲线的数目。应当指出，对任意的环面曲面 $X$，Mikhalkin 公式并不计算 Gromov--Witten 不变量，因为总会有来自非环面横截的曲线（例如有分量映入 $X$ 的环面边界的曲线）对这些 Gromov--Witten 不变量的贡献。然而对 $\PP^2$，我们确实得到 Gromov--Witten 不变量。这使得可以纯粹用热带方式定义 $\PP^2$ 的量子上同调，并由此得到完整的 A-模型半无限 Hodge 结构变分。

这就引出如下问题：能否用热带方式描述 B-模型，使得 A-模型与 B-模型一致这一点变得透明？乍看之下，如何用热带方式计算势函数 $W_0$ 的扰动以及涉及这些函数的振荡积分，是完全不清楚的。然而，正如本章将要展示的，这些计算确实可以用热带方式进行，特别是可以给出 B-模型的完全热带描述。本章完全基于论文 [42]。

% [Tutor Note: 开篇的 $\langle T_2^{3d-1}\rangle_{0,d}$ 与定义 5.16 中的记号一致：$T_i=T_1^i$ 是 $\PP^2$ 的量子上同调基（见第 2 章）。核对锚点：$d=1$ 时 $\langle T_2^{2}\rangle_{0,1}=1$（过 2 个点的直线数），而例 5.17 的 $\langle\psi^{3d-2}T_2\rangle^{\mathrm{trop}}_{0,1}=1$，两者都是经典计数；一般 $d$ 的数值见例 5.17 与定理 5.18。]

\subsection{扰动后的 Landau--Ginzburg 势}

我们的主要目标是给出 $W_0$ 及其扰动的一个典范描述，即给出 $W_0$ 的“正确”的万有展开。重要的线索是 Cho 与 Oh [16] 把势 $W_0$ 解释为 Maslov 指标 2 的全纯圆盘。这里不在辛几何的层面解释这一点，而是直接转到热带版本。同时，我们先在任意环面簇这一更大的框架中工作，从而把第 2 章已讨论过的 $\PP^n$ 的镜像的描述，放到 Givental 引入的环面簇镜像对称的更一般框架中（见 [33]）。

\subsubsection{Givental 的环面簇镜像}

先一次性固定如下记号：格 $M=\Z^n$，对偶格 $N=\operatorname{Hom}_\Z(M,\Z)$，$M_{\R}=M\otimes_\Z\R$，$N_{\R}=N\otimes_\Z\R$。固定 $M_{\R}$ 中的一个完全扇 $\Sigma$。我们假设 $X_\Sigma$ 是非异环面簇。沿用 §1.3 的记号：令 $T_\Sigma$ 是由 $\Sigma^{[1]}$ 的元素（即 $\Sigma$ 的一维锥）生成的自由 Abel 群，其中 $\rho\in\Sigma^{[1]}$ 对应生成元 $t_\rho$。有映射 $r:T_\Sigma\to M$，由 $r(t_\rho)=m_\rho$（$\rho$ 的本原生成元）给出；非异性的假设蕴含 $r$ 是满射。于是有自然正合列
\[
  0\longrightarrow K_\Sigma\longrightarrow T_\Sigma\xrightarrow{\ r\ }M\longrightarrow 0
\]

% [原文 p.174]
它定义了 $K_\Sigma$。把该正合列对偶化，得到
\[
  0\to N\to \operatorname{Hom}_\Z(T_\Sigma,\Z)\to\operatorname{Pic}X_\Sigma\to 0,
\]
正如 §3.1.2 中所讨论的那样。再与 $\C^\times$ 作张量积，得到正合列
\[
  0\longrightarrow N\otimes\C^\times\longrightarrow \operatorname{Hom}(T_\Sigma,\C^\times)
  \xrightarrow{\ \kappa\ }\operatorname{Pic}X_\Sigma\otimes\C^\times\longrightarrow 0.
\]
应当把 $\kappa$ 理解为给出 $X_\Sigma$ 的镜像族。我们定义 $X_\Sigma$ 的 \emph{Kähler 模空间}为
\[
  \cM_\Sigma:=\operatorname{Pic}X_\Sigma\otimes\C^\times=\operatorname{Spec}\C[K_\Sigma],
\]
于是有态射
\[
  \kappa:\operatorname{Hom}(T_\Sigma,\C^\times)=\operatorname{Spec}\C[T_\Sigma]\to\cM_\Sigma.
\]
注意 $\kappa$ 在 $\cM_\Sigma$ 的闭点上的纤维典范同构于 $\operatorname{Spec}\C[M]$。

与第 2 章一致，在过渡到复流形之后，可以过渡到 $\cM_\Sigma$ 的万有覆盖，我们把它写成向量空间
\[
  \tilde\cM_\Sigma:=\operatorname{Pic}X_\Sigma\otimes\C,
\]
其中的映射 $\tilde\cM_\Sigma\to\cM_\Sigma$ 由 $D\otimes y\mapsto D\otimes e^y$ 给出。然后令
\[
  \check\cX_\Sigma:=\operatorname{Hom}(T_\Sigma,\C^\times)\times_{\cM_\Sigma}\tilde\cM_\Sigma,
\]
即把 $\kappa$ 给出的族拉回到 $\tilde\cM_\Sigma$ 上。到此为止，我们还没有引入该模空间的形式加厚，一般也不会这样做。

我们取 $W_0$ 为 $\check\cX_\Sigma$ 上由下式定义的函数
\begin{equation}
  W_0:=\sum_{\rho\in\Sigma^{[1]}}z^{t_\rho}. \tag{5.1}
\end{equation}
把它看作 $\operatorname{Spec}\C[T_\Sigma]$ 上的函数是有意义的，从而也是 $\check\cX_\Sigma$ 上的函数。我们把 $(\check\cX_\Sigma,W_0)$ 看作一族 Landau--Ginzburg 模型，它们是 $X_\Sigma$ 的镜像。

\begin{example}
取 $\Sigma$ 为 $\PP^n$ 的扇，则这恰好给出 §2.2.3 中定义的 $\PP^n$ 的镜像族 $\check\cX\to\tilde\cM$，其中形式坐标为 $t_0=t_2=\cdots=t_n=0$。
\end{example}

% [Tutor Note: (5.1) 中的 $z^{t_\rho}$ 是 §1.3 中记号 $z^m$（$m\in T_\Sigma$ 对应的单项式）的特例；对 $\PP^n$，$\Sigma^{[1]}=\{\rho_0,\rho_1,\rho_2\}$，$z^{t_i}$ 就是例 5.9 中的 $x_i$。核对锚点：$\kappa$ 的纤维是 $\operatorname{Spec}\C[M]$，故 $z^{t_\rho}$ 在纤维上的限制就是通常的坐标函数。]

\subsubsection{热带圆盘与环面簇的镜像}

下面限制到曲面情形，取 $M=\Z^2$。

我们在第 1 章已经定义了热带曲线的概念。这里我们定义热带圆盘的概念（在 [85] 中它被称为带停点的热带曲线）。

\begin{definition}
设 $\overline{\Gamma}$ 是如 §1.3 中那样的带权、连通、无二价顶点的有限图，并附加选取一个一价顶点 $V_{\mathrm{out}}$，它邻接于唯一一条边 $E_{\mathrm{out}}$。令
\[
  \Gamma':=\bigl(\overline{\Gamma}\setminus\overline{\Gamma}^{[0]}_{\infty}\bigr)\cup\{V_{\mathrm{out}}\}\subseteq\overline{\Gamma}.
\]
再设 $\Gamma'$ 的第一 Betti 数为零（即 $\Gamma'$ 是一棵树，带一条紧的外部边与若干条非紧的外部边）。于是 $M_{\R}$ 中以 $\Gamma'$ 为定义域的 \emph{参数化 $d$ 点热带圆盘}指的是一个包含映射 $\{p_1,\ldots,p_d\}\hookrightarrow\Gamma^{[1]}_{\infty}\setminus\{E_{\mathrm{out}}\}$，记作 $p_i\mapsto E_{p_i}$，使得 $w(E)=0$ 当且仅当

% [原文 p.175]
$E=E_{p_i}$ 对某个 $i$ 成立，以及一个连续映射 $h:\Gamma'\to M_{\R}$，它满足与定义 1.11 相同的条件，只是在 $V_{\mathrm{out}}$ 处\emph{没有}平衡条件。
\end{definition}

热带圆盘的同构
\[
  h_1:(\Gamma_1',p_1,\ldots,p_d)\to M_{\R}
  \quad\text{与}\quad
  h_2:(\Gamma_2',p_1,\ldots,p_d)\to M_{\R}
\]
指的是一个同胚 $\Phi:\Gamma_1'\to\Gamma_2'$，它保持标记边与权，且满足 $h_1=h_2\circ\Phi$。一个\emph{热带圆盘}就是参数化热带圆盘的同构类。

热带圆盘 $h:(\Gamma',p_1,\ldots,p_d)\to M_{\R}$ 的\emph{组合型}定义为 $\overline{\Gamma}$ 连同标记点、权、$V_{\mathrm{out}}$ 的同胚类，并对每个顶点 $V$ 与含 $V$ 的每条边 $E$，还附带 $h(E)$ 在 $M$ 中沿离开 $V$ 方向的本原切向量。

上述定义在任何维数都说得通，但限制到二维的意义在于：我们可以使用 Mikhalkin 对重数的定义（定义 1.21），对圆盘采用同样的定义，但不计 $V_{\mathrm{out}}$：

\begin{definition}
设 $\operatorname{rank}M=2$。设 $h:\Gamma'\to M_{\R}$ 是带标记热带圆盘，且 $\overline{\Gamma}$ 只有价为 1 和 3 的顶点。则圆盘 $h$ 的\emph{重数}为
\[
  \operatorname{Mult}(h):=\prod_{V\in\Gamma^{[0]}}\operatorname{Mult}_V(h),
\]
其中 $\operatorname{Mult}_V(h)$ 如定义 1.21 所定义。特别地，$V_{\mathrm{out}}$ 不产生贡献。
\end{definition}

% [Tutor Note: 定义 5.3 说 $V_{\mathrm{out}}$ 不产生贡献；注意 §5.1.2 的记号中 $\overline{\Gamma}$ 是原图、$\Gamma'$ 是剪去全部一价顶点 $\overline{\Gamma}^{[0]}_{\infty}$ 后再加回 $V_{\mathrm{out}}$ 所得的树，乘积 $\prod_{V\in\Gamma^{[0]}}$ 应理解为只对重数非平凡的顶点取值（原文记号在此处有轻微不一致，照译未改）。核对锚点：条件 $w(E)=0$ 当且仅当 $E=E_{p_i}$（$i=1,\ldots,d$）。]

现在固定 $M\cong\Z^2$ 以及 $M_{\R}$ 中的完全有理多面体扇 $\Sigma$，由此按 §5.1.1 得到 $T_\Sigma$ 与 $r:T_\Sigma\to M$。

\begin{definition}
称圆盘 $h$ 是 $X_\Sigma$ 中的\emph{热带圆盘}，如果对每条 $E\in\Gamma^{[1]}_{\infty}\setminus\{E_{\mathrm{out}}\}$，$h(E)$ 要么是一个点，要么是某个 $\rho\in\Sigma^{[1]}$ 的平移。
\end{definition}

类似于定义 1.12，若该圆盘有 $d_\rho$ 条无界边是 $\rho\in\Sigma^{[1]}$ 的平移（按权计数），则 $h$ 的\emph{次数}为
\[
  \Delta(h):=\sum_{\rho\in\Sigma^{[1]}}d_\rho t_\rho\in T_\Sigma.
\]

固定一般位置的点 $P_1,\ldots,P_k\in M_{\R}$，并固定一个一般位置的\emph{基点} $Q\in M_{\R}$。下文谈到一般位置的点时，指的是存在一个开稠密子集（通常是若干个余维数至少为 1 的多面体之并的补集），使得 $(P_1,\ldots,P_k,Q)\in M_{\R}^{k+1}$ 落在该开子集中。这个开子集的选取将依赖于具体需要。

把点 $P_1,\ldots,P_k$ 与环
\[
  R_k:=\frac{\C[u_1,\ldots,u_k]}{(u_1^2,\ldots,u_k^2)}
\]
中的变量 $u_1,\ldots,u_k$ 相联系。

\begin{definition}
设 $h:(\Gamma',p_1,\ldots,p_d)\to M_{\R}$ 是 $X_\Sigma$ 中的热带圆盘，且 $h(V_{\mathrm{out}})=Q$，$h(p_j)=P_{i_j}$，$1\le i_1<\cdots<i_d\le k$。（这一排序消除了标记点标号的 $d!$ 重歧义。）我们说 $h$ 是以 $Q$ 为边界的 $(X_\Sigma,P_1,\ldots,P_k)$ 中的热带圆盘。
\end{definition}

% [原文 p.176]
圆盘 $h$ 的 \emph{Maslov 指标}为
\[
  \operatorname{MI}(h):=2\bigl(|\Delta(h)|-d\bigr).
\]

“Maslov 指标”一词在边界含于辛流形的拉格朗日子流形中的伪全纯圆盘的语境下有确切定义。特别地，Maslov 指标出现在这类圆盘族的虚拟维数的公式中。我们给出上面的定义，是为了让一个完全类似的公式成立，它由下面的引理给出。

\begin{lemma}
若 $P_1,\ldots,P_k,Q$ 取自一般位置，则 $(X_\Sigma,P_1,\ldots,P_k)$ 中以 $Q$ 为边界的 Maslov 指标 $2n$ 的热带圆盘之集是 $(n-1)$ 维多面体复形。边界任意的 Maslov 指标 $2n$ 的热带圆盘之集是 $(n+1)$ 维多面体复形。
\end{lemma}

\para{证明。}这正是引理 1.20 中所作的那种标准的热带维数计数论证。这里略述其论证。

固定带 $d$ 个标记点、次数为 $\Delta$ 的热带圆盘的一个组合型。若该组合型是一般的，则定义域 $\Gamma'$ 除 $V_{\mathrm{out}}$ 外只有三价顶点。这样的树有 $|\Delta|+d$ 条无界边，从而有 $|\Delta|+d-1$ 条有界边（含 $E_{\mathrm{out}}$）。于是给定组合型的热带圆盘 $h:\Gamma'\to M_{\R}$ 完全由 $h(V_{\mathrm{out}})\in M_{\R}$ 的位置以及各条有界边的仿射长度确定。这在所有 $d$ 点、次数为 $\Delta$ 的热带圆盘的模空间 $\cM^{\mathrm{disk}}_{\Delta,d}(X_\Sigma)$ 中产生一个胞腔。该胞腔的闭包为 $(\R_{\ge0})^{|\Delta|+d-1}\times M_{\R}$。此外，给定次数的圆盘只有有限多个组合型。因此 $\cM^{\mathrm{disk}}_{\Delta,d}(X_\Sigma)$ 是有限维的 $(|\Delta|+d+1)$ 维多面体复形。再者，有分片线性映射 $\mathrm{ev}:\cM^{\mathrm{disk}}_{\Delta,d}(X_\Sigma)\to M_{\R}^d$，把圆盘 $h$ 映为多元组 $(h(p_1),\ldots,h(p_d))$。令 $E\subseteq\cM^{\mathrm{disk}}_{\Delta,d}(X_\Sigma)$ 为在 $\mathrm{ev}$ 下映入 $M_{\R}^d$ 中余维数 $\ge1$ 的胞腔的那些胞腔之并；于是 $h(E)$ 是 $M_{\R}^d$ 的闭子集。因此，若 $(P_{i_1},\ldots,P_{i_d})\in M_{\R}^d$（$1\le i_1<\cdots<i_d\le k$ 为互不相同的指标）不在该闭子集中，则 $\mathrm{ev}^{-1}(P_{i_1},\ldots,P_{i_d})$ 是 $\cM^{\mathrm{disk}}_{\Delta,d}(X_\Sigma)$ 的余维数 $2d$ 的子集。于是给定次数 $\Delta$、边界任意的热带圆盘在 $(X_\Sigma,P_1,\ldots,P_k)$ 中的模空间维数为 $|\Delta|+1-d=\operatorname{MI}(h)/2+1$。类似地，若再固定一般位置的边界点 $Q$，则维数为 $\operatorname{MI}(h)/2-1$，如所断言。

\begin{definition}
给定一般位置的 $P_1,\ldots,P_k,Q\in M_{\R}$，设
\[
  h:(\Gamma',p_1,\ldots,p_d)\to M_{\R}
\]
是 $(X_\Sigma,P_1,\ldots,P_k)$ 中以 $Q$ 为边界的 Maslov 指标 2 的带标记热带圆盘。则可以把 $\C[T_\Sigma]\otimes_\C R_k[[y_0]]$ 中的一个单项式联系到 $h$：
\[
  \operatorname{Mono}(h):=\operatorname{Mult}(h)z^{\Delta(h)}u_{I(h)},
\]
其中 $z^{\Delta(h)}\in\C[T_\Sigma]$ 是对应于 $\Delta(h)\in T_\Sigma$ 的单项式，子集 $I(h)\subseteq\{1,\ldots,k\}$ 由
\[
  I(h):=\{i\mid h(p_j)=P_i\text{ 对某个 }j\}
\]
定义，且
\[
  u_{I(h)}=\prod_{i\in I(h)}u_i.
\]
\end{definition}

% [原文 p.177]
定义 \emph{$k$ 点 Landau--Ginzburg 势}为
\[
  W_k(Q):=y_0+\sum_h\operatorname{Mono}(h)\in\C[T_\Sigma]\otimes_\C R_k[[y_0]],
\]
其中求和取遍 $(X_\Sigma,P_1,\ldots,P_k)$ 中以 $Q$ 为边界的全部 Maslov 指标 2 的圆盘 $h$。由引理 5.6，对一般的 $P_1,\ldots,P_k,Q$，这是一个有限和。

我们定义 §5.1.1 中所定义的 Kähler 模空间 $\tilde\cM_\Sigma$ 的 \emph{$k$ 阶加厚}为环化空间
\[
  \tilde\cM_{\Sigma,k}:=(\tilde\cM_\Sigma,\cO_{\tilde\cM_{\Sigma,k}}),
\]
其中对 $U\subseteq\tilde\cM_\Sigma$，$\cO_{\tilde\cM_{\Sigma,k}}(U)$ 的元素是形如
\[
  \sum_{n=0}^{\infty}\ \sum_{I\subseteq\{1,\ldots,k\}} f_{n,I}y_0^nu_I
\]
的表达式，其中对每个 $n$ 与 $I$，$f_{n,I}$ 是 $U$ 上的全纯函数。

类似地，我们以完全相同的方式定义加厚后的镜像族
\[
  \check\cX_{\Sigma,k}:=(\check\cX_\Sigma,\cO_{\check\cX_{\Sigma,k}}).
\]
于是有族
\[
  \kappa:\check\cX_{\Sigma,k}\to\tilde\cM_{\Sigma,k}.
\]
由构造，$W_k(Q)$ 是 $\check\cX_{\Sigma,k}$ 上的正则函数，故可以把它看作给出一族 Landau--Ginzburg 势。

相对微分层 $\Omega^1_{\check\cX_{\Sigma,k}/\tilde\cM_{\Sigma,k}}$ 典范同构于平凡局部自由层 $M\otimes_\Z\cO_{\check\cX_{\Sigma,k}}$，其中 $m\otimes1$ 对应于微分
\begin{equation}
  d\log m:=\frac{d(z^{\overline{m}})}{z^{\overline{m}}}; \tag{5.2}
\end{equation}
这里 $\overline{m}\in T_\Sigma$ 是 $m\in M$ 的任意提升，且 $d\log m$ 作为相对微分与提升的选取无关。于是 $\bigwedge^2M\cong\Z$ 的一个生成元的选取决定了处处非零的相对全纯二形式 $\Omega$，它在相差一个符号的意义下典范。确切地说，若 $e_1,e_2\in M$ 是正定向的基，则
\begin{equation}
  \Omega=d\log e_1\wedge d\log e_2. \tag{5.3}
\end{equation}

% [Tutor Note: (5.2) 中“与提升无关”的依据：若把 $\overline m$ 换成 $\overline m+k$（$k\in K_\Sigma$），则 $z^{\overline m}$ 被乘以 $z^k$，而 $z^k$ 沿 $\check\cX_{\Sigma,k}\to\tilde\cM_{\Sigma,k}$ 的纤维是常数（因为 $K_\Sigma$ 是 $r$ 的核，$z^k$ 来自底空间），故 $d\log z^k$ 作为\emph{相对}微分为零。核对锚点：(5.3) 只用到 $\bigwedge^2M$ 的生成元，即 $\Omega$ 与 $\check\cX_{\Sigma,k}$ 的具体坐标无关。]

\begin{remark}
在 $M_{\R}\cong\R^2$ 中给定扇 $\Sigma$，若取 $k=0$，则以 $Q$ 为边界的 Maslov 指标 2 的圆盘恰是形如 $Q+\rho$（$\rho\in\Sigma^{[1]}$）的圆盘。这些圆盘的重数全为 1，因为它们没有顶点。于是
\[
  W_0(Q)=y_0+\sum_{\rho\in\Sigma^{[1]}}z^{t_\rho}.
\]
除 $y_0$ 这一项之外（它在后文讨论中的作用仅在于重现 $J$-函数的正确行为），这正是 (5.1)。对熟悉 Cho 与 Oh [16] 工作的读者，我们指出：这些热带圆盘恰好与那里所分类的全纯圆盘重合，因此我们用热带方式重现了 [16] 中给出的 Landau--Ginzburg 势的描述。
\end{remark}

% [原文 p.178]
\begin{figure}[htbp]
\centering
\includegraphics[width=0.5\textwidth]{f5_1.png}
\figtitle{1}{$k=1$ 时额外的那一个 Maslov 指标 2 热带圆盘。}
\end{figure}

函数 $W_k(Q)$ 于是意在成为 $W_0(Q)$ 的“正确”扰动，其意义是：$W_k(Q)$ 中出现的参数与平坦坐标密切相关。然而对一般的 $\Sigma$，这并不成立，主要问题在于 $X_\Sigma$ 的环面边界中会出现一些 $\PP^1$ 的副本，它们不能形变为与 $X_\Sigma$ 的大环面轨道相交的曲线。这是热带几何中的标准问题：热带几何“看不见”这些曲线。只要 $X_\Sigma$ 是射影空间的乘积，这就不是问题；因此特别地，我们下面限制到 $X_\Sigma=\PP^2$ 的情形。这就回到第 2 章中详细研究过的例子。

\begin{example}
取 $\Sigma$ 为 $\PP^2$ 的扇（第 1 章的图 14 中已画出），故 $X_\Sigma=\PP^2$。这里 $T_\Sigma=\Z^3$，基 $t_0,t_1,t_2$ 分别对应 $\rho_0,\rho_1,\rho_2$，我们记 $x_i$ 为单项式 $z^{t_i}\in\C[T_\Sigma]$。此时 $\tilde\cM_\Sigma=\C$，设其坐标为 $y_1$，于是空间 $\tilde\cM_{\Sigma,k}$ 的底层拓扑空间是 $\C$。镜像族 $\check\cX_{\Sigma,k}$ 由 $\C^3\times\tilde\cM_{\Sigma,k}$ 中的方程 $x_0x_1x_2=e^{y_1}$ 定义，而映射
\[
  \kappa:\check\cX_{\Sigma,k}\to\tilde\cM_{\Sigma,k}
\]
则由投影给出，即
\[
  \kappa^*(y_0)=y_0,\qquad
  \kappa^*(y_1)=y_1,\qquad
  \kappa^*(u_i)=u_i.
\]
如注 5.8 所讨论，
\[
  W_0(Q)=y_0+x_0+x_1+x_2.
\]
若取 $k=1$，即在 $\PP^2$ 中标记一个点，则得到额外的那个圆盘，如图 1 所示；若取 $k=2$ 且 $P_1,P_2$ 如图 2 选取，则得到三个额外的圆盘。注意势依赖于点 $P_1,\ldots,P_k$ 以及 $Q$ 的具体选取。在所述的例子中，分别有
\[
  \begin{aligned}
    W_1(Q)&=y_0+x_0+x_1+x_2+u_1x_1x_2\\
    W_2(Q)&=y_0+x_0+x_1+x_2+u_1x_0x_1+u_2x_0x_1+u_1u_2x_0x_1^2.
  \end{aligned}
\]
\end{example}

% [原文 p.179]
\begin{figure}[htbp]
\centering
\includegraphics[width=0.86\textwidth]{f5_2.png}
\figtitle{2}{$k=2$ 时额外的 Maslov 指标 2 热带圆盘。}
\end{figure}

\subsection{热带后代不变量}

正如我们从第 2 章所知道的，$\PP^2$ 的镜像对称用振荡积分给出了 $\PP^2$ 的 $J$-函数的一个表达式。$J$-函数又涉及引力后代 Gromov--Witten 不变量。另一方面，本章的主旨是表明振荡积分可以用某些热带描述的对象自然地求值。特别地，我们需要定义引力后代不变量的热带版本，而我们将对 $\PP^2$ 的情形来做这件事。

穿过 $3d-1$ 个点的 $d$ 次有理曲线的热带计数，最终是由“热带曲线对应于全纯曲线”这一实际证明所推动的：Mikhalkin 在 [80] 中的原始证明所得的计数，与第 4 章中计数的产生方式非常相似。然而对引力后代而言，要说明这些公式的动机则多少更困难一些。并没有类似于第 4 章论证的直接证明来说明我们定义的后代不变量确实与全纯版本一致。我们只能这样来说明这些定义的动机：这些定义正是通过计算振荡积分而得到的。本章的主要结果表明：对热带后代不变量所给的定义是正确的，当且仅当 $\PP^2$ 的镜像对称（定理 2.44）成立。

也可以尝试在纯热带框架中处理这个问题：考虑热带曲线的模空间，把 $\psi$-类定义为这些模空间上的热带循环，并应用热带相交理论。

% [原文 p.180]
这一途径被 Markwig 与 Rau 在 [77] 中采用，但至今尚不能重新得到我们这里给出的全部公式。

下面继续取 $M=\Z^2$，$\Sigma$ 是 $M_{\R}$ 中定义射影环面簇 $X_\Sigma$ 的扇。

\begin{definition}
设 $P_1,\ldots,P_k\in M_{\R}$ 一般，$S\subseteq M_{\R}$ 是子集。定义
\[
  \cM^{\mathrm{trop}}_{\Delta,n}(X_\Sigma,P_1,\ldots,P_k,\psi^\nu S)
\]
为 $X_\Sigma$ 中次数为 $\Delta$ 的有理 $(n+1)$ 点热带曲线
\[
  h:(\Gamma,p_1,\ldots,p_n,x)\to M_{\R}
\]
的模空间，其中
\begin{enumerate}[label=(\arabic*),leftmargin=2.2em,itemsep=1pt]
\item $h(p_j)=P_{i_j}$，$1\le i_1<\cdots<i_n\le k$。
\item 边 $E_x$ 附着于 $\Gamma$ 的一个顶点 $V_x$；该顶点的价记为 $\operatorname{Val}(x)$。则
  \[
    \operatorname{Val}(x)=\nu+3.
  \]
\item $h(x)\in S$。
\item $\Gamma$ 的每条无界边的权为 $0$ 或 $1$。
\end{enumerate}
\end{definition}

\begin{lemma}
对一般的 $P_1,\ldots,P_k\in M_{\R}$，
\begin{enumerate}[label=(\arabic*),leftmargin=2.2em,itemsep=1pt]
\item $\cM^{\mathrm{trop}}_{\Delta,n}(X_\Sigma,P_1,\ldots,P_k,\psi^\nu M_{\R})$ 是维数为 $|\Delta|-n-\nu$ 的多面体复形。
\item 当 $C$ 是 $M_{\R}$ 中一条热带曲线的一般平移时，$\cM^{\mathrm{trop}}_{\Delta,n}(X_\Sigma,P_1,\ldots,P_k,\psi^\nu C)$ 是维数为 $|\Delta|-n-\nu-1$ 的多面体复形。
\item 当 $Q\in M_{\R}$ 是一般点时，$\cM^{\mathrm{trop}}_{\Delta,n}(X_\Sigma,P_1,\ldots,P_k,\psi^\nu Q)$ 是维数为 $|\Delta|-n-\nu-2$ 的多面体复形。
\end{enumerate}
\end{lemma}

\para{证明。}这很直接，与引理 5.6 相同。维数计数如下。固定曲线的一般组合型，使得 $\Gamma$ 除与 $E_x$ 相邻的顶点外全为三价顶点，而该顶点是 $(\nu+3)$ 价的。这样的树有 $|\Delta|+n+1$ 条无界边，从而有 $|\Delta|+n+1-(\nu+3)$ 条有界边。这种组合型的曲线于是由 $h(x)\in M_{\R}$ 的位置以及各条有界边的长度确定，给出形如 $(\R_{\ge0})^{|\Delta|+n-\nu-2}\times M_{\R}$ 的胞腔。固定 $h(p_1),\ldots,h(p_n)$ 后，就得到 (1) 中模空间所需的维数
\[
  |\Delta|+n-\nu-2n.
\]
这给出 (1)。至于 (2) 与 (3)，考虑映射
\[
  \mathrm{ev}_x:\cM^{\mathrm{trop}}_{\Delta,n}(X_\Sigma,P_1,\ldots,P_k,\psi^\nu M_{\R})\to M_{\R},
\]
由 $\mathrm{ev}_x(h):=h(x)$ 给出。令 $E_1\subseteq\cM^{\mathrm{trop}}_{\Delta,n}(X_\Sigma,P_1,\ldots,P_k,\psi^\nu M_{\R})$ 为映入 $M_{\R}$ 中余维数 $\ge1$ 的集合的那些胞腔之并，$E_2$ 为映入 $M_{\R}$ 中的点的那些胞胞之并。于是需要选取平移 $C$，使得 $C\cap h(E_1)$ 是零维的且 $C\cap h(E_2)=\varnothing$。类似地，需要选取 $Q\notin h(E_1)$。此时 $\mathrm{ev}_x^{-1}(C)$ 或 $\mathrm{ev}_x^{-1}(Q)$ 就是 (2)、(3) 中所需的模空间，且具有所需的维数。

% [Tutor Note: 引理 5.11 的维数计数：固定 $h(p_1),\ldots,h(p_n)$ 应减去 $2n$ 而不是 $n$，因为每个条件 $h(p_j)=P_j$ 在二维的 $M_{\R}$ 上是 2 个方程；于是 (1) 的维数 $|\Delta|+n-\nu-2n=|\Delta|-n-\nu$，与引理陈述一致。核对锚点：胞腔维数 $(\R_{\ge0})^{|\Delta|+n-\nu-2}\times M_{\R}$ 为 $|\Delta|+n-\nu$。]

我们可以通过在选定的顶点处把曲线切开，把亏格零的热带曲线 $h:\Gamma\to M_{\R}$ 描述为一族圆盘。下面的引理对于后文用圆盘描述这些曲线至关重要。

% [原文 p.181]
\begin{lemma}
设 $P_1,\ldots,P_k\in M_{\R}$ 一般，$S\subseteq M_{\R}$ 是子集。设
\[
  h\in\cM^{\mathrm{trop}}_{\Delta,n}(X_\Sigma,P_1,\ldots,P_k,\psi^\nu S).
\]
用 $\Gamma_1',\ldots,\Gamma_{\nu+2}'$ 记 $\Gamma\setminus E_x$ 的连通分支的闭包，$h_i:\Gamma_i'\to M_{\R}$ 为 $h$ 的限制。每个圆盘 $h_i$ 被视为以那些满足 $E_p\subseteq\Gamma_i'$ 的点 $p\in\{p_1,\ldots,p_n\}$ 为标记。这里有一个特殊情形需要考虑：若 $E_x$ 与 $E_{p_i}$ 有公共顶点 $V$，则我们也把边 $E_{p_i}$ 从考虑中舍弃，于是得到圆盘 $h_1,\ldots,h_{\nu+1}$。（注意，由于 $i\ne j$ 时 $h(p_i)\ne h(p_j)$，$E_{p_i}$ 与 $E_{p_j}$ 绝不会共有公共顶点。）
\begin{enumerate}[label=(\arabic*),leftmargin=2.2em,itemsep=1pt]
\item 若 $S=M_{\R}$ 且 $n=|\Delta|-\nu$，则或者
  \begin{enumerate}[label=(\alph*),leftmargin=2.2em,itemsep=1pt]
  \item $E_x$ 不与任何 $E_{p_i}$ 共有顶点，此时除两个 $i$ 之外对所有的 $i$ 都有 $\operatorname{MI}(h_i)=2$，而对这两个 $i$ 有 $\operatorname{MI}(h_i)=0$；
  \item $E_x$ 与某个 $E_{p_i}$ 共有顶点，此时对所有的 $i$ 都有 $\operatorname{MI}(h_i)=2$。
  \end{enumerate}
\item 若 $S=C$ 是 $M_{\R}$ 中一条热带曲线的一般平移，且 $n=|\Delta|-\nu-1$，则除一个 $i$ 之外对所有的 $i$ 都有 $\operatorname{MI}(h_i)=2$，而对该 $i$ 有 $\operatorname{MI}(h_i)=0$。
\item 若 $S=\{Q\}$ 其中 $Q$ 是一般点，且 $n=|\Delta|-\nu-2$，则对所有的 $i$ 都有 $\operatorname{MI}(h_i)=2$。
\end{enumerate}
\end{lemma}

\para{证明。}首先注意，对 $n$ 的条件以及 $P_1,\ldots,P_k$、$C$、$Q$ 的一般性，由前一引理保证所考虑的模空间是零维的。若某个圆盘 $h_i$ 能在保持其边界 $h_i(x)$ 固定的情况下形变，则这给出 $h$ 的非平凡形变，而它不存在。于是由引理 5.6，在所有情形下都必须有 $\operatorname{MI}(h_i)\le2$。令 $n_i$ 为 $h_i$ 上标记点的数目。注意
\[
  \begin{aligned}
    \sum_i\frac{\operatorname{MI}(h_i)}{2}
    &=\sum_i\bigl(|\Delta(h_i)|-n_i\bigr)\\
    &=\begin{cases}
      |\Delta(h)|-(n-1) & \text{情形 (1)(b)}\\
      |\Delta(h)|-n & \text{其他情形}
    \end{cases}\\
    &=\begin{cases}
      \nu & \text{情形 (1)(a)}\\
      \nu+1 & \text{情形 (1)(b)}\\
      \nu+1 & \text{情形 (2)}\\
      \nu+2 & \text{情形 (3)}
    \end{cases}
  \end{aligned}
\]
由于除情形 (1)(b) 外都有 $\nu+2$ 个圆盘（情形 (1)(b) 有 $\nu+1$ 个圆盘），结论随之成立。

现在可以定义出现在 Givental $J$-函数中的后代 Gromov--Witten 不变量的热带类比。本章自此只考虑 $X_\Sigma=\PP^2$ 的情形，其中 $\Sigma$ 是由
$m_0=(-1,-1)$、$m_1=(1,0)$、$m_2=(0,1)$ 生成的射线构成的扇，$t_0,t_1,t_2$ 是 $T_\Sigma$ 的生成元，且 $r(t_i)=m_i$。令
\[
  \Delta_d:=d(t_0+t_1+t_2)\in T_\Sigma;
\]
次数为 $\Delta_d$ 的曲线通常应看作 $\PP^2$ 中 $d$ 次热带曲线。

\begin{definition}
固定一般位置的点 $Q,P_1,P_2,\cdots\in M_{\R}$。设 $L$ 是以 $Q$ 为顶点的热带直线（$\PP^2$ 的扇中一维锥之并的平移）。

% [原文 p.182]
对 $\PP^2$ 中带标记点 $x$ 的热带曲线 $h$，令 $n_0(x)$、$n_1(x)$、$n_2(x)$ 分别为与 $E_x$ 共有公共顶点、且在 $h$ 下的像落在方向 $m_0$、$m_1$、$m_2$ 上的无界射线的数目。如引理 5.12 那样，用 $h_1,\ldots$ 记从 $\Gamma$ 中删去 $E_x$ 所得的热带圆盘，$h_i$ 的出边为 $E_{i,\mathrm{out}}$。令 $m(h_i)=w(E_{i,\mathrm{out}})m^{\mathrm{prim}}(h_i)$，其中 $m^{\mathrm{prim}}(h_i)\in M$ 是切于 $h_i(E_{i,\mathrm{out}})$、指向离开 $h(x)$ 方向的本原向量。

定义
\begin{align*}
  \operatorname{Mult}^0_x(h)&=\frac{1}{n_0(x)!\,n_1(x)!\,n_2(x)!}\\
  \operatorname{Mult}^1_x(h)&=
  -\frac{\sum_{k=1}^{n_0(x)}\frac{1}{k}+\sum_{k=1}^{n_1(x)}\frac{1}{k}+\sum_{k=1}^{n_2(x)}\frac{1}{k}}
  {n_0(x)!\,n_1(x)!\,n_2(x)!}\\
  \operatorname{Mult}^2_x(h)&=
  \frac{\Bigl(\sum_{l=0}^{2}\sum_{k=1}^{n_l(x)}\frac{1}{k}\Bigr)^{2}
  +\sum_{l=0}^{2}\sum_{k=1}^{n_l(x)}\frac{1}{k^{2}}}
  {2\,n_0(x)!\,n_1(x)!\,n_2(x)!}
\end{align*}

% [Tutor Note: 三式中的 $\operatorname{Mult}^1_x$ 带负号，$\operatorname{Mult}^2_x$ 的分子是“先对 $l,k$ 求和、再整体平方”加上“平方和”，两者不可混同。核对锚点：取例 5.17 的 $n_0=1,n_1=0,n_2=2$，得 $\operatorname{Mult}^0_x=1/2$、$\operatorname{Mult}^1_x=-(1+1+1/2)/(1!0!2!)=-5/4$、$\operatorname{Mult}^2_x=((1+1+1/2)^2+(1+1+1/4))/(2\cdot1!0!2!)=17/8$，与 p.186 的计算一致。]

\begin{enumerate}[label=(\arabic*),leftmargin=2.2em,itemsep=2pt]
\item 定义
  \[
    \bigl\langle P_1,\ldots,P_{3d-2-\nu},\psi^\nu Q\bigr\rangle^{\mathrm{trop}}_{0,d}
  \]
  为
  \[
    \sum_h\operatorname{Mult}(h),
  \]
  其中求和取遍全部带标记热带有理曲线
  \[
    h\in\cM^{\mathrm{trop}}_{\Delta_d,3d-2-\nu}(P_1,\ldots,P_{3d-2-\nu},\psi^\nu Q).
  \]
  并定义
  \[
    \operatorname{Mult}(h):=\operatorname{Mult}^0_x(h)\prod_{\substack{V\in\Gamma^{[0]}\\ V\notin E_x}}\operatorname{Mult}_V(h).
  \]
\item 定义
  \[
    \bigl\langle P_1,\ldots,P_{3d-1-\nu},\psi^\nu L\bigr\rangle^{\mathrm{trop}}_{0,d}
  \]
  为和
  \[
    \sum_h\operatorname{Mult}(h),
  \]
  其中求和取遍全部带标记热带有理曲线
  \[
    h:(\Gamma,p_1,\ldots,p_{3d-1-\nu},x)\to M_{\R},
  \]
  满足 $h(p_i)=P_i$，且满足下面两个条件之一。
  \begin{enumerate}[label=(\alph*),leftmargin=2.2em,itemsep=2pt]
  \item $h\in\cM^{\mathrm{trop}}_{\Delta_d,3d-1-\nu}(P_1,\ldots,P_{3d-1-\nu},\psi^\nu L)$。
    此外，$\Gamma$ 中除 $E_x$ 外与 $E_x$ 共有公共顶点的任何无界边都不映入 $L\setminus\{Q\}$ 的含 $h(x)$ 的连通分支。由引理 5.12，恰有一个 $j$（$1\le j\le\nu+2$）使 $\operatorname{MI}(h_j)=0$。再设 $L\setminus\{Q\}$ 的含 $h(x)$ 的连通分支为 $Q+\R_{>0}m_i$。则定义
    \[
      \operatorname{Mult}(h)=|m(h_j)\wedge m_i|\operatorname{Mult}^0_x(h)
      \prod_{\substack{V\in\Gamma^{[0]}\\ V\notin E_x}}\operatorname{Mult}_V(h).
    \]
    这里 $m(h_j)\wedge m_i\in\bigwedge^2M\cong\Z$，故绝对值有意义。

% [原文 p.183]
  \item $\nu\ge1$ 且
    \[
      h\in\cM^{\mathrm{trop}}_{\Delta_d,3d-1-\nu}(P_1,\ldots,P_{3d-1-\nu},\psi^{\nu-1}Q).
    \]
    则
    \[
      \operatorname{Mult}(h)=\operatorname{Mult}^1_x(h)
      \prod_{\substack{V\in\Gamma^{[0]}\\ V\notin E_x}}\operatorname{Mult}_V(h).
    \]
  \end{enumerate}
\item 定义
  \[
    \bigl\langle P_1,\ldots,P_{3d-\nu},\psi^\nu M_{\R}\bigr\rangle^{\mathrm{trop}}_{0,d}
  \]
  为和
  \[
    \sum_h\operatorname{Mult}(h),
  \]
  其中求和取遍全部带标记热带有理曲线
  \[
    h:(\Gamma,p_1,\ldots,p_{3d-\nu},x)\to M_{\R},
  \]
  满足 $h(p_i)=P_i$，且满足下列之一。
  \begin{enumerate}[label=(\alph*),leftmargin=2.2em,itemsep=2pt]
  \item
    \[
      h\in\cM^{\mathrm{trop}}_{\Delta_d,3d-\nu}(P_1,\ldots,P_{3d-\nu},\psi^\nu M_{\R}),
    \]
    且 $E_x$ 不与任何 $E_{p_i}$ 共有顶点。此外，$\Gamma$ 中除 $E_x$ 外与 $E_x$ 共有公共顶点的任何无界边都不映入 $M_{\R}\setminus L$ 的含 $h(x)$ 的连通分支。由引理 5.12，恰有两个不同的 $j_1,j_2$（$1\le j_1,j_2\le\nu+2$）使 $\operatorname{MI}(h_{j_i})=0$。则定义
    \[
      \operatorname{Mult}(h)=|m(h_{j_1})\wedge m(h_{j_2})|\operatorname{Mult}^0_x(h)
      \prod_{\substack{V\in\Gamma^{[0]}\\ V\notin E_x}}\operatorname{Mult}_V(h).
    \]
  \item
    \[
      h\in\cM^{\mathrm{trop}}_{\Delta_d,3d-\nu}(P_1,\ldots,P_{3d-\nu},\psi^\nu M_{\R}),
    \]
    且 $E_x$ 与某条 $E_{p_i}$ 共有顶点。此外，$\Gamma$ 中除 $E_x$ 与 $E_{p_i}$ 外与 $E_x$ 共有公共顶点的任何无界边都不映入 $M_{\R}\setminus L$ 的含 $h(x)$ 的连通分支。则定义
    \[
      \operatorname{Mult}(h)=\operatorname{Mult}^0_x(h)
      \prod_{\substack{V\in\Gamma^{[0]}\\ V\notin E_x}}\operatorname{Mult}_V(h).
    \]
  \item $\nu\ge1$ 且
    \[
      h\in\cM^{\mathrm{trop}}_{\Delta_d,3d-\nu}(P_1,\ldots,P_{3d-\nu},\psi^{\nu-1}L).
    \]
    此外，$\Gamma$ 中除 $E_x$ 外与 $E_x$ 共有公共顶点的任何无界边都不映入 $L\setminus\{Q\}$ 的含 $h(x)$ 的连通分支。由引理 5.12，恰有一个 $j$（$1\le j\le\nu+1$）使 $\operatorname{MI}(h_j)=0$。设 $L\setminus\{Q\}$ 的含 $h(x)$ 的连通分支为 $Q+\R_{>0}m_i$。则定义
    \[
      \operatorname{Mult}(h)=|m(h_j)\wedge m_i|\operatorname{Mult}^1_x(h)
      \prod_{\substack{V\in\Gamma^{[0]}\\ V\notin E_x}}\operatorname{Mult}_V(h).
    \]

% [原文 p.184]
  \item $\nu\ge2$ 且
    \[
      h\in\cM^{\mathrm{trop}}_{\Delta_d,3d-\nu}(P_1,\ldots,P_{3d-\nu},\psi^{\nu-2}Q).
    \]
    则
    \[
      \operatorname{Mult}(h)=\operatorname{Mult}^2_x(h)
      \prod_{\substack{V\in\Gamma^{[0]}\\ V\notin E_x}}\operatorname{Mult}_V(h).
    \]
  \end{enumerate}
\end{enumerate}
对所有情形 $S=\{Q\}$、$L$ 或 $M_{\R}$，对 $\sigma\in\Sigma$ 定义
\begin{equation}
  \bigl\langle P_1,\ldots,P_{3d-\nu-(2-\dim S)},\psi^\nu S\bigr\rangle^{\mathrm{trop}}_{d,\sigma} \tag{5.4}
\end{equation}
为 $\bigl\langle P_1,\ldots,P_{3d-\nu-(2-\dim S)},\psi^\nu S\bigr\rangle^{\mathrm{trop}}_{0,d}$ 中来自那些满足 $h(x)$ 落在 $\sigma+Q$ 的内部的曲线 $h$ 的贡献。在 (1) 中，唯一的贡献来自 $\sigma=\{0\}$；在 (2) 中，贡献来自 $\Sigma$ 的零维与一维锥；在 (3) 中，贡献来自 $\Sigma$ 的全部锥。
\end{definition}

% [Tutor Note: 定义 5.13 (2)(a) 与 (3)(a)(c) 中“除 $E_x$（以及 (3)(b) 中的 $E_{p_i}$）外，与 $E_x$ 有公共顶点的无界边都不映入某连通分支”这一条件很关键：它保证 $x$ 所在圆盘的出边方向唯一确定 $m_i$，从而 $|m(h_j)\wedge m_i|$ 有定义。核对锚点：$\PP^2$ 的扇有三条射线，故 $L\setminus\{Q\}$ 与 $M_{\R}\setminus L$ 各恰有 3 个连通分支，分别对应 $m_0,m_1,m_2$。]

\begin{remark}
(1) 注意对一般的 $Q,P_1,\ldots$，所涉及的模空间全是零维的，故这些和有意义。

(2) 定义 5.13 中 (1) 的公式在 $\nu=0$ 时给出穿过 $3d-1$ 个点的 $d$ 次有理曲线数目的热带曲线计数公式，即 Mikhalkin 公式（定理 4.4）。当 $\nu>0$ 时，它与 Markwig 与 Rau 在 [77] 中给出的公式一致。特别地，由该文的结果，
\[
  \bigl\langle P_1,\ldots,P_{3d-2-\nu},\psi^\nu Q\bigr\rangle^{\mathrm{trop}}_{0,d}
  =\bigl\langle T_2^{3d-2-\nu},\psi^\nu T_2\bigr\rangle_{0,d}.
\]
我们给出的其余公式则更为神秘，除了本章所给的镜像对称论证之外，尚不知有任何依据。

(3) 容易看出，$\bigl\langle P_1,\ldots,P_{3d-1},\psi^0L\bigr\rangle^{\mathrm{trop}}_{0,d}$ 等于穿过 $3d-1$ 个点的有理曲线数目的 $d$ 倍。事实上，对该数目的唯一贡献来自定义 5.13 的 (2)(a)。对每条带 $3d-1$ 个标记点、穿过 $P_1,\ldots,P_{3d-1}$ 的热带有理曲线 $h:\Gamma\to M_{\R}$，对 $h^{-1}(L)$ 的每个点用 $x$ 标记该点，就得到一个贡献。此情形下重数中的因子 $|m(h_j)\wedge m_i|\operatorname{Mult}^0_x(h)$ 给出 $h(\Gamma)$ 与 $L$ 在 $h^{-1}(L)$ 的每个点处的相交重数，如例 1.6 所定义。由热带 Bézout 定理，来自 $h$ 的总贡献于是为 $(h(\Gamma)\cdot L)\operatorname{Mult}(h)=d\operatorname{Mult}(h)$。因此
\[
  \bigl\langle P_1,\ldots,P_{3d-1},\psi^0L\bigr\rangle^{\mathrm{trop}}_{0,d}
  =d\bigl\langle P_1,\ldots,P_{3d-1}\bigr\rangle^{\mathrm{trop}}_{0,d}
  =d\bigl\langle T_2^{3d-1}\bigr\rangle_{0,d}
  =\bigl\langle T_2^{3d-1},T_1\bigr\rangle_{0,d},
\]
最后一步由 Mikhalkin 公式与除子公理得到。

(4) $\bigl\langle P_1,\ldots,P_{3d},\psi^0M_{\R}\bigr\rangle^{\mathrm{trop}}_{0,d}=0$。事实上，可能的贡献只来自定义 5.13 的 (3)(a)，但穿过 $3d$ 个一般点的 $d$ 次有理曲线并不存在。于是
\[
  \bigl\langle P_1,\ldots,P_{3d},\psi^0M_{\R}\bigr\rangle^{\mathrm{trop}}_{0,d}
  =\bigl\langle T_2^{3d},T_0\bigr\rangle_{0,d},
\]
两者都为零，后者由基本类公理得到。
\end{remark}

我们将在 §5.5 中证明：

\begin{theorem}
定义 5.13 中定义的不变量与诸 $P_i$ 以及 $Q$ 的选取无关。
\end{theorem}

% [Tutor Note: 定理 5.15 的证明在 §5.5（本块之外）；在此之前出现的所有 $\sum_h$ 都必须理解为“对零维模空间求和”，这一点由注 5.14(1) 与引理 5.11 保证。核对锚点：定义 5.16 中的三条等式右端都含有 $Q$、$L$、$M_{\R}$，定理 5.15 正是使这些记号有意义的前提。]

这使我们能作如下定义。

% [原文 p.185]
\begin{definition}
定义
\[
  \begin{aligned}
    \bigl\langle T_2^{3d-2-\nu},\psi^\nu T_2\bigr\rangle^{\mathrm{trop}}_{0,d}
    &:=\bigl\langle P_1,\ldots,P_{3d-2-\nu},\psi^\nu Q\bigr\rangle^{\mathrm{trop}}_{0,d}\\
    \bigl\langle T_2^{3d-1-\nu},\psi^\nu T_1\bigr\rangle^{\mathrm{trop}}_{0,d}
    &:=\bigl\langle P_1,\ldots,P_{3d-1-\nu},\psi^\nu L\bigr\rangle^{\mathrm{trop}}_{0,d}\\
    \bigl\langle T_2^{3d-\nu},\psi^\nu T_0\bigr\rangle^{\mathrm{trop}}_{0,d}
    &:=\bigl\langle P_1,\ldots,P_{3d-\nu},\psi^\nu M_{\R}\bigr\rangle^{\mathrm{trop}}_{0,d}
  \end{aligned}
\]
其中诸 $P_i$ 与 $Q$ 已取为一般位置。

若 $m+i+\nu\ne3d$，我们定义
\[
  \bigl\langle T_2^{m},\psi^\nu T_i\bigr\rangle^{\mathrm{trop}}_{0,d}:=0.
\]
类比于例 2.30，我们定义 $\PP^2$ 的\emph{热带 $J$-函数}
\[
  \begin{aligned}
    J^{\mathrm{trop}}_{\PP^2}
    &:=e^{(y_0T_0+y_1T_1)/\hbar}\cup\Biggl(
    T_0+\sum_{i=0}^{2}\Biggl(\hbar^{-1}y_2\delta_{2,i}\\
    &\qquad+\sum_{d\ge1}\sum_{\nu\ge0}
    \bigl\langle T_2^{3d+i-2-\nu},\psi^\nu T_{2-i}\bigr\rangle^{\mathrm{trop}}_{0,d}
    \hbar^{-(\nu+2)}e^{dy_1}\frac{y_2^{3d+i-2-\nu}}{(3d+i-2-\nu)!}\Biggr)T_i\Biggr)\\
    &=:\sum_{i=0}^{2}J^{\mathrm{trop}}_iT_i.
  \end{aligned}
\]
\end{definition}

\begin{example}
热带后代不变量（相对而言）容易计算。例如考虑 $\bigl\langle\psi^{3d-2}T_2\bigr\rangle^{\mathrm{trop}}_{0,d}$。在 $Q$ 处有价 $3d+1$ 的顶点（含标记边）的 $d$ 次热带曲线只有一条，即分别在 $(-1,-1)$、$(1,0)$、$(0,1)$ 三个方向上各有 $d$ 条权为 1 的leg的曲线，它因此贡献重数 $1/(d!)^3$，故
\[
  \bigl\langle\psi^{3d-2}T_2\bigr\rangle^{\mathrm{trop}}_{0,d}=\frac{1}{(d!)^{3}}.
\]
其次考虑 $\bigl\langle T_2,\psi^{3d-3}T_2\bigr\rangle^{\mathrm{trop}}_{0,d}$。固定某点 $P_1\in M_{\R}$，注意任何对此不变量有贡献的曲线都在 $Q$ 处有 $3d$ 价顶点（含标记边），并分解为 $3d-1$ 个以 $Q$ 为边界的 Maslov 指标 2 圆盘，其中恰有一个经过点 $P_1$。取比如图 1 中所示的 $P_1$，则对此不变量有贡献的唯一曲线即图 3 所示。于是
\[
  \bigl\langle T_2,\psi^{3d-3}T_2\bigr\rangle^{\mathrm{trop}}_{0,d}
  =\frac{1}{d!\,(d-1)!\,(d-1)!}.
\]
两者都给出正确的非热带后代不变量。这将由定理 5.18、推论 5.19 与定理 2.44 推出。

我们再给出若干热带后代 Gromov--Witten 不变量贡献的例子。在图 4 中，我们看到一条通过把四个以 $Q$ 为边界的 Maslov 指标 2 圆盘粘合而得到的曲线：竖直线旁附的数目 2 表示把该竖直无界边取两次，而水平线段旁附的 2 告诉我们该边的权为 2。除 $Q$ 外所有顶点对重数的贡献为 $2\times2$。$Q$ 处的顶点是 5 价的（记住该顶点同时也是边 $E_x$ 的端点），

% [原文 p.186]
\begin{figure}[htbp]
\centering
\includegraphics[width=0.56\textwidth]{f5_3.png}
\figtitle{3}{对 $\bigl\langle P_1,\psi^{3d-3}Q\bigr\rangle^{\mathrm{trop}}_{0,d}$ 有贡献的那一条曲线。数目 $d$ 与 $d-1$ 表示的不是边的权，而是该边的副本数。}
\end{figure}

其中 $n_0(x)=1$、$n_1(x)=0$、$n_2(x)=2$。于是
\[
  \begin{aligned}
    \operatorname{Mult}^0_x(h)&=1/(1!\,0!\,2!)=1/2,\\
    \operatorname{Mult}^1_x(h)&=-(1+1+1/2)/(1!\,0!\,2!)=-5/4,\\
    \operatorname{Mult}^2_x(h)&=\bigl((1+1+1/2)^{2}+1+1+1/4\bigr)/(2\cdot1!\,0!\,2!)=17/8.
  \end{aligned}
\]
于是该曲线对 $\bigl\langle P_1,\ldots,P_5,\psi^2Q\bigr\rangle^{\mathrm{trop}}_{0,3}$ 贡献 $4/2=2$，对
$\bigl\langle P_1,\ldots,P_5,\psi^3L\bigr\rangle^{\mathrm{trop}}_{0,3}$ 贡献 $4\times(-5/4)=-5$，对
$\bigl\langle P_1,\ldots,P_5,\psi^4M_{\R}\bigr\rangle^{\mathrm{trop}}_{0,3}$ 贡献 $4\times(17/8)=17/2$。

图 5 中我们考虑另一个例子，其中虚线为 $L$。此情形下，删去 $E_x$ 把曲线分解为若干圆盘后可见，其中唯一 Maslov 指标为零的圆盘是过 $P_6$ 的对角线。于是因子 $|m(h_j)\wedge m_1|$ 为 1。我们有 $n_0(x)=n_1(x)=0$、$n_2(x)=2$。（注意，由于标记边映射到 $P_6$，过 $P_6$ 的对角射线不是与 $E_x$ 共有公共顶点的无界边的像！）于是 $\operatorname{Mult}^0_x(h)=1/2$，$\operatorname{Mult}^1_x(h)=-3/4$。因此该曲线对 $\bigl\langle P_1,\ldots,P_6,\psi^2L\bigr\rangle^{\mathrm{trop}}_{0,3}$ 的贡献为 $2$，对 $\bigl\langle P_1,\ldots,P_6,\psi^3M_{\R}\bigr\rangle^{\mathrm{trop}}_{0,3}$ 的贡献为 $-3$。

图 6 画出一条 4 次曲线，其映入 $h(x)$ 的顶点价为 7，且 $n_0(x)=n_1(x)=0$（必然如此，否则该曲线不会有贡献），$n_2(x)=4$。于是 $\operatorname{Mult}^0_x(h)=1/24$，而分解该曲线所得的两个 Maslov 指标零圆盘，其出边切向量分别为 $(-1,-2)$ 与 $(1,-2)$。于是因子 $|m(h_{j_1})\wedge m(h_{j_2})|$ 为 4。所有其他顶点的重数均为 1，只有两个顶点的重数为 2，故该曲线对 $\bigl\langle P_1,\ldots,P_8,\psi^4M_{\R}\bigr\rangle^{\mathrm{trop}}_{0,4}$ 的总贡献为 $2/3$。

我们留给读者给出一个对热带后代不变量定义中 (3)(b) 型后代不变量有贡献的曲线的例子。
\end{example}

% [Tutor Note: 例 5.17 第二式的 $1/(d!(d-1)!(d-1)!)$ 可用 Mult 公式核对：该曲线在 $Q$ 处为 $3d$ 价顶点（含标记边），按图 3 三条leg上的副本数 $d-1,d-1,d$ 分配，得 $n_0=n_1=d-1$、$n_2=d$，故 $\operatorname{Mult}^0_x(h)=1/(d!(d-1)!(d-1)!)$，其余顶点重数为 1。注意图 3 的说明强调 $d,d-1$ 是副本数而非权。]

% [原文 p.187]
\begin{figure}[htbp]
\centering
\includegraphics[width=0.86\textwidth]{f5_4.png}
\figtitle{4}{}
\end{figure}

\begin{figure}[htbp]
\centering
\includegraphics[width=0.86\textwidth]{f5_5.png}
\figtitle{5}{}
\end{figure}

\subsection{主 B-模型陈述}

本节继续取 $M=\Z^2$、$\Sigma$ 为 $M_{\R}$ 中定义 $\PP^2$ 的扇。于是对 $k\ge0$，由定义 5.7 给出
\[
  \kappa:\check\cX_{\Sigma,k}\to\tilde\cM_{\Sigma,k}.
\]
给定点 $P_1,\ldots,P_k,Q\in M_{\R}$ 的一般选取，便得到 Landau--Ginzburg 势 $W_k(Q)$。注意模去 $u_1,\ldots,u_k$ 后，$W_k(Q)=W_0(Q)$。空间 $\check\cX_{\Sigma,k}$ 带有由 (5.3) 定义的相对全纯二形式。

设 $R$ 是 $\tilde\cM_{\Sigma,k}\times\C^\times$ 上的局部系统，其在 $(u,\hbar)$ 上的纤维为
\[
  H_2\bigl(\kappa^{-1}(u),\Re(W_0(Q)/\hbar)\ll0\bigr).
\]
这个局部系统与结构层是否加厚无关：它只承载底层拓扑空间 $\tilde\cM_{\Sigma,k}\times\C^\times$（即 $\tilde\cM_\Sigma\times\C^\times$）上的纯拓扑信息。

% [Tutor Note: 本块末句按 STYLE_GUIDE §10「跨页句子开始者写完」补完：原文该句在 p.187 末起头（“it carries purely topological information on the underlying topological space of”），在 p.188 接上 “$\tilde\cM_{\Sigma,k}\times\C^\times$, which is $\tilde\cM_\Sigma\times\C^\times$.”。下一块 c05b 应从 “In fact, this local system coincides with the one given in §2.2.3 in the case of $\PP^2$.” 起译。]

% [Tutor Note: 本块涉及的两处编号引用需与相邻块核对：定理 5.18、推论 5.19 在 §5.5（印刷 p.209 之后）；§2.2.3 与 (2.38) 见第 2 章；定理 4.4、例 1.6、定义 1.11/1.12/1.21、引理 1.20 均按原文编号。]

